\section{Explicit stable arclength and the original optical weight}
\label{sec:eo-optical}

The optical comparison in Input~\ref{input:v112-geometric-pressure}
has a direct proof with explicit constants.  We give that proof in
the true collision coordinates, including the sign of the physical
derivative and the change to the original optical coordinates.
It completes the two optical inequalities in that input.  It does
not prove its separate whole-cylinder pressure bound.
The qualitative comparison and its proof in historical module 238
remain part of the retained source; the contribution here is a
complete quantitative norm and cocycle verification.

Throughout this section the particle has unit speed, the scatterers
are the original radius-$R$ disks on the triangular lattice, and
$R\in[9/20,47/100]$.  Only true regular first-collision trajectories
are considered.  Put
\begin{equation}\label{eq:eo-constants}
 a=\frac{100}{47},\qquad K=\frac{170}{9},\qquad H=19,\qquad
 \epsilon_{\rm ang}=\frac{2a}{H^2},\qquad
 \tau_-=\frac3{50},\qquad
 \Lambda=1+2\tau_-a=\frac{59}{47}.
\end{equation}
Here $a\le\varkappa=R^{-1}\le20/9$ and
$K=20/9+\tau_-^{-1}$; $1+K^2<H^2$.
The constants $a,K,\epsilon_{\rm ang}$ are cone constants, unrelated
to the original section probability $c_*$ or the incidence guard
$q_0$.  Write $s=R\alpha$, $p=\sin\varphi$, and
$c=\cos\varphi>0$.  Norms on tangent vectors in $(s,\varphi)$
are Euclidean norms, with no adapted metric.
The original probability is still
$\nu=(4\pi R)^{-1}c\,ds\,d\varphi=d\alpha\,dp/(4\pi)$,
and the section-conditioned probability is its original
$c_*^{-1}$ restriction.  The radius is fixed along each
collision cocycle, so this coordinate conversion introduces
no factor depending on the collision count.

\begin{lemma}[The physical differential and two cone directions]
\label{lem:eo-differential}
At a regular flight from $x$ to $x'=T_Rx$, write its physical
length as $\tau$, its incidences as $c,c'$, and its endpoint
curvatures as $\varkappa,\varkappa'$.  With the counterclockwise
boundary tangent and the above signed outgoing momentum,
\begin{equation}\label{eq:eo-one-flight}
 D T_R(x)=-M(x),\qquad
 M(x)=\frac1{c'}
 \begin{pmatrix}
   \tau\varkappa+c&\tau\\
   \tau\varkappa\varkappa'+\varkappa c'
          +\varkappa'c&\tau\varkappa'+c'
 \end{pmatrix},\qquad
 \det M(x)=\frac c{c'}.
\end{equation}
On the set $M_R'$ regular in both infinite time directions there
are invariant unstable and stable lines with slopes
\[
             v^u(x)\in[a,K],\qquad v^s(x)\in[-K,-a].
\]
They are measurable.  If
$e(x)=|\sin\angle(E^s(x),E^u(x))|$, then
\begin{equation}\label{eq:eo-angle}
                    \epsilon_{\rm ang}\le e(x)\le1.
\end{equation}
The positive arclength cocycle $j_m(x)=J^sT_R^m(x)$ satisfies
\begin{equation}\label{eq:eo-stable-decay}
             j_m(x)\le\frac{H}{\epsilon_{\rm ang}}\Lambda^{-m}
                         \quad(m\ge0).
\end{equation}
\end{lemma}

\begin{proof}
Let $n(\alpha)=(\cos\alpha,\sin\alpha)$,
$t(\alpha)=(-\sin\alpha,\cos\alpha)$, and let $v$ be the
unit direction of the actual flight.  At departure and arrival,
respectively,
\[
                  v=cn+pt,\qquad v=-c'n'+p't'.
\]
For a variation of the departure state,
$\delta v=v^\perp(\delta\varphi+\varkappa\delta s)$.
Differentiating $Q'-Q=\tau v$ and taking its scalar product with
$v^\perp$ gives
\[
     -c'\delta s'-c\delta s
                   =\tau(\delta\varphi+\varkappa\delta s).
\]
The incoming direction angle at the arrival point is
$\alpha'+\pi-\varphi'$; thus
$\delta\varphi'=\varkappa'\delta s'
                         -\varkappa\delta s-\delta\varphi$.
These two identities prove \eqref{eq:eo-one-flight}, including
its minus sign.  Expanding the determinant of the displayed
numerator leaves exactly $cc'$.

Scalar signs do not affect slopes.  The map induced by $M(x)$ on
an input slope $v\ge a$ is
\begin{equation}\label{eq:eo-slope-map}
 f_x(v)=\varkappa'
     +\frac{c'(\varkappa+v)}{\tau(\varkappa+v)+c}.
\end{equation}
Consequently $a\le f_x(v)\le\varkappa'+\tau^{-1}\le K$.
On $[a,K]$ its derivative satisfies
\begin{equation}\label{eq:eo-slope-contraction}
 0<f_x'(v)=
     \frac{cc'}{[\tau(\varkappa+v)+c]^2}
 \le\frac1{4\tau(\varkappa+v)}
 \le\frac1{8\tau_-a}=\frac{47}{48}<1.
\end{equation}
The first inequality uses $c'\le1$ and
$(b+c)^2\ge4bc$ with $b=\tau(\varkappa+v)>0$.
The images at $x$ of the cone interval $[a,K]$ transported from
$T_R^{-n}x$ are nested nonempty compact intervals, with diameter
at most $(K-a)(47/48)^n$.  Their intersection is therefore one
line $E^u(x)$, and shifting the sequence proves its invariance.
Every finite approximation is measurable on the countable regular
collision partition; its limit is measurable.

Time reversal in these coordinates is
$I(s,\varphi)=(s,-\varphi)$ and $T_R^{-1}=IT_RI$ on the
regular domain.  Applying the same construction to reversed
flights gives the invariant line $E^s(x)$ with negative slope.
For unit vectors on these two lines,
\[
 e(x)=\frac{v^u(x)-v^s(x)}
          {\sqrt{1+v^u(x)^2}\sqrt{1+v^s(x)^2}}
                  \ge\frac{2a}{H^2}.
\]
This proves \eqref{eq:eo-angle}.

For completeness these invariant cone lines have the claimed
stable character in the actual Euclidean norm.  An unstable
vector with $s$-component $\delta s_i$ obeys
\[
 c_{i+1}|\delta s_{i+1}|
  =[\tau_i(\varkappa_i+v_i^u)+c_i]|\delta s_i|
  \ge\Lambda c_i|\delta s_i|,
\]
since $c_i\le1$.  Hence, writing
$u_m(x)=J^uT_R^m(x)$,
\begin{equation}\label{eq:eo-unstable-growth}
               u_m(x)\ge
                 \frac{c_0}{Hc_m}\Lambda^m.
\end{equation}
Taking absolute areas of the images of the two unit invariant
vectors, and multiplying the exact one-step determinants, gives
\begin{equation}\label{eq:eo-area-cocycle}
        u_m(x)j_m(x)\frac{e(T_R^mx)}{e(x)}
                         =\frac{c_0}{c_m}.
\end{equation}
Combining \eqref{eq:eo-unstable-growth} with this identity yields
\eqref{eq:eo-stable-decay}; its $m=0$ case also holds because
$j_0=1$.  Reversal gives the corresponding backward contraction
on $E^u$.  Thus these are the true stable and unstable lines, not
auxiliary test directions.  No ergodic, pressure or mixing theorem
has entered their construction.
\end{proof}

For $m\ge1$ put
\[
       \mathsf D_m(x)=M(T_R^{m-1}x)\cdots M(x),\qquad
       b_m(x)=(\mathsf D_m(x))_{12}>0.
\]
This matrix product is distinct from the scalar physical-word
partition sum $P_m(R)$ in module 266.

\begin{lemma}[One comparison of columns, without accumulated constants]
\label{lem:eo-columns}
For every $m\ge1$ and $x\in M_R'$,
\begin{equation}\label{eq:eo-column-bound}
       \frac{2a}{H}\,b_m(x)\le u_m(x)
                               \le2KH\,b_m(x).
\end{equation}
The original inverse endpoint derivative is exactly
\begin{equation}\label{eq:eo-original-omega}
 \omega(x)=|\partial_{p_0}\alpha_m|^{-1}
                   =\frac{Rc_0}{b_m(x)}
                   =\frac{Rc_0c_m}{\mathfrak b_m(x)}.
\end{equation}
Here $\mathfrak b_m$ is the original Jacobi optical denominator
of \eqref{eq:v109-optical-product}.
\end{lemma}

\begin{proof}
For the two columns $C_1,C_2$ of one matrix $M(x)$, each component
of $C_1/C_2$ lies in $[a,K]$.  In the first row the ratio is
$\varkappa+c/\tau$; in the second it is
$\varkappa+\varkappa'c/(\tau\varkappa'+c')$.
Left multiplication by a positive matrix preserves the
componentwise inequalities $aC_2\le C_1\le KC_2$.
They therefore hold for the columns of the whole product, with
exactly the same constants.

The second column of the whole product has slope in $[a,K]$.
For the first factor this follows from
$(\tau\varkappa'+c')/\tau$; later factors preserve the cone.
Consequently $b_m\le|C_2|\le Hb_m$.  A unit unstable vector at
the source is, up to sign, $(1,v)/\sqrt{1+v^2}$ with
$v\in[a,K]$.  Its image is, up to the common sign $(-1)^m$,
$(C_1+vC_2)/\sqrt{1+v^2}$.  Positivity and the column inequalities
give
\[
 \frac{2a}{H}b_m
 \le\frac{a+v}{\sqrt{1+v^2}}|C_2|
 \le u_m
 \le\frac{K+v}{\sqrt{1+v^2}}|C_2|
 \le2KHb_m.
\]
Only the initial column comparison and the terminal cone bound
have been used; no constant is raised to the collision count.

With $\alpha_0$ fixed, the coordinate conversions are
$\delta\varphi_0=\delta p_0/c_0$ and
$\delta\alpha_m=\delta s_m/R$.  Thus
$|\partial_{p_0}\alpha_m|=b_m/(Rc_0)$.
There is one factor $R$, irrespective of $m$.
The agreement with the retained denominator can also be checked
as an exact conjugacy.  Set
\[
 B_i=\begin{pmatrix}c_i&0\\\varkappa_i&1\end{pmatrix},
 \qquad C_i=\begin{pmatrix}1&0\\2\varkappa_i/c_i&1\end{pmatrix},
 \qquad S_i=\begin{pmatrix}1&\tau_i\\0&1\end{pmatrix}.
\]
Direct multiplication gives
\begin{equation}\label{eq:eo-optical-conjugacy}
                    B_{i+1}M_iB_i^{-1}=C_{i+1}S_i.
\end{equation}
In the product, all intermediate $B_i$ cancel.
The top row of $B_m^{-1}C_m$ is $(1/c_m,0)$ and the second
column of $B_0$ is $(0,1)^{\mathsf T}$.  Hence
$b_m=\mathfrak b_m/c_m$, with precisely the internal optical
factors of \eqref{eq:v109-optical-product}.  For $m=1$ this says
$b_1=\tau_0/c_1$ and $\mathfrak b_1=\tau_0$, so the empty
internal product is included.
\end{proof}

\begin{theorem}[Explicit original-coordinate optical comparison]
\label{thm:eo-optical}
On $M_R'$, for every $R\in[9/20,47/100]$ and every $m\ge1$,
\begin{equation}\label{eq:eo-final}
 \frac1{30000}\,c_m(x)j_m(x)
       \le\omega(x)\le30000\,c_m(x)j_m(x),\qquad
                   \frac{j_1(x)}{c_0(x)}\le6400.
\end{equation}
The same constants work for every original exact return,
displacement and collision label, every physical roof interval,
and all original initial and terminal section tests.
In fact this comparison requires no positive endpoint cap beyond
regularity.  It includes trajectories arbitrarily close to
interior or endpoint grazing.
\end{theorem}

\begin{proof}
Equations \eqref{eq:eo-area-cocycle} and
\eqref{eq:eo-original-omega} imply the exact ratio
\begin{equation}\label{eq:eo-ratio}
 \frac{\omega(x)}{c_m(x)j_m(x)}
             =R\,\frac{u_m(x)}{b_m(x)}
                          \frac{e(T_R^mx)}{e(x)}.
\end{equation}
The previous two lemmas therefore place this ratio between
$4R_-a^2/H^3$ and $R_+KH^3/a$.
All constants are fixed numbers, and direct rational arithmetic
gives
\begin{equation}\label{eq:eo-rational-budget}
 \frac{H^3}{4R_-a^2}<30000,\qquad
 \frac{R_+KH^3}{a}<30000,\qquad
 \frac{H^3}{4\tau_-a^2}<6400.
\end{equation}
These prove the first two inequalities in \eqref{eq:eo-final}.

For one flight, the $s$-component of the image of a unit
unstable vector and its terminal norm give
\[
 c_1u_1
  =[\tau_0(\varkappa_0+v^u_0)+c_0]\,
       \frac{\sqrt{1+(v^u_1)^2}}{\sqrt{1+(v^u_0)^2}}
                    \ge\frac{2\tau_-a}{H}.
\]
The one-step area identity now gives
\[
 \frac{j_1}{c_0}
    =\frac{e(x)}{c_1u_1e(T_Rx)}
    \le\frac{H}{2\tau_-a\,\epsilon_{\rm ang}}
    =\frac{H^3}{4\tau_-a^2}<6400.
\]
There is no inverse incidence left over.  None of these
calculations discards an original event test: the inequalities
hold before any such test is inserted.
\end{proof}

\begin{corollary}[Precisely the remaining pressure premise]
\label{cor:eo-pressure-scope}
The two optical assertions in
Input~\ref{input:v112-geometric-pressure} hold with
$C_0=30000$ and $C_1=6400$.  Its remaining premise is
\[
              \sup_{R,l}
                 \sum_{w\in\mathcal W_l(R)}
                    \sup_{x\in[w]_R\cap M_R'}j_l(x)\le Q_*,
                 \qquad P_0(R)=1.
\]
Under this separately retained whole-cylinder premise, all
consequences of module 266 may take
\[
 C_\star=\frac{30000Q_*}{2c_*\lambda_\circ},
                    \qquad K_\star=6400D Q_*.
\]
The invariant-line restriction does not remove positive
finite-record collision mass.
\end{corollary}

\begin{proof}
The theorem supplies exactly the two optical estimates, in the
same norm convention on every nonempty stable block.  On an
empty block the cocycle remains $j_0=1$; the optical inverse
derivative is used only for records with at least one flight and
is not defined or invoked for $m=0$.
The finite marked-word factorization in module 266 therefore
uses these constants without any change to its indices or empty
blocks.  It still requires its displayed pressure sum.

For each fixed radius the one-step singularities and grazing
states form a finite union of regular arcs and their endpoints
in each lifted collision chart.  They have collision area zero.
On every regular branch the determinant in
\eqref{eq:eo-one-flight} is nonzero.  Countable regular
coordinate charts and change of variables show that every
forward and backward iterate of a one-step null singular set is
null.  Removing their countable union thus leaves $M_R'$ with
full collision measure.  Restriction to the original finite
event and pushforward by the original roof preserve equality
of these positive bulk measures.  Equivalently the local
endpoint-coordinate map has its nonzero derivative in
\eqref{eq:eo-original-omega}, so the corresponding
two-dimensional area formulas are unchanged almost everywhere.
No common null set for uncountably many radii is required by
$\sup_R\operatorname*{ess\,sup}_t$.
This bulk statement assigns no zero trace to a one-sided
singular boundary or to a paired endpoint measure.
\end{proof}
